\documentclass[10pt,a4paper]{amsart}
\usepackage[margin=19mm]{geometry}
\usepackage{amsmath,amssymb,mathtools}
\usepackage[scr=rsfs,cal=euler]{mathalfa}
\usepackage{xfrac}
\usepackage[unicode,hidelinks,bookmarks=false]{hyperref}
\newcommand{\ie}{i.e.\ }
\newcommand{\cf}{cf.\ }
\newcommand{\resp}{respectively\ }
\newcommand{\Subsection}{\subsection}
\providecommand{\nobreakdash}{\nobreak\hbox{-}}
\setlength{\emergencystretch}{2em}
\multlinegap0pt
\usepackage{bm}
\usepackage{bbm}
\usepackage{dsfont}
\usepackage{xspace}
\usepackage{cancel}
\usepackage{mathtools}
\usepackage{xcolor}
\usepackage{amsthm}
\makeatletter
\@ifundefined{theorem}{\newtheorem{theorem}{Theorem}}{}
\@ifundefined{corollary}{\newtheorem{corollary}[theorem]{Corollary}}{}
\@ifundefined{lemma}{\newtheorem{lemma}[theorem]{Lemma}}{}
\@ifundefined{proposition}{\newtheorem{proposition}[theorem]{Proposition}}{}
\makeatother
\newcommand{\bb}[1]{{\mathbb #1}}
\newcommand{\mc}[1]{{\mathcal #1}}
\newcommand{\mf}[1]{{\mathfrak #1}}
\newcommand{\ms}[1]{{\mathscr #1}}
\newcommand{\mtt}[1]{{\mathtt #1}}
\title[From one-dimensional diffusion processes metastable behaviou]{From one-dimensional diffusion processes metastable behaviour to parabolic equations asymptotics}
\author{Claudio Landim,, Christian Maura}
\date{Source: arXiv:2505.20217v1 (CC BY 4.0)}
\begin{document}
\maketitle

\noindent\emph{Source and licence.} From one-dimensional diffusion processes metastable behaviour to parabolic equations asymptotics. Version arXiv:2505.20217v1 (CC BY 4.0), distributed under the Creative Commons Attribution 4.0 International licence, as recorded in the arXiv licence metadata for this exact version and shown on the abstract page https://arxiv.org/abs/2505.20217v1. Full rights notice: licenses/arxiv-2505.20217-CC-BY-4.0.txt.\par\smallskip

\noindent\textit{Editorial scope.} Counted unit: the Lemma whose source label is \texttt{l36} in the article (source0.tex lines 4488--4502, source label \texttt{l36}), delivered together with the article's own complete original proof of it (source0.tex lines 4504--4621, one \texttt{proof} environment of 118 lines). No mathematics is written, repaired or re-derived: statement and proof are the article's own text line for line. Required context retained uncounted: 71 (lines 1585--1588); l07 (lines 3977--3986); l11 (lines 4088--4104); 33 (lines 4286--4291). Every definition, display and label of those lines is retained unchanged. Omitted from this package and not redistributed: 1--1584, 1589--3976, 3987--4087, 4105--4285, 4292--4487, 4503--4503, 4622--6481. Nothing inside a retained excerpt is omitted or altered: every retained body is delivered exactly as the article's lines are written, so no omission receipt of any kind accompanies this package. Citations: the retained excerpts carry no citation command at all, so no bibliography entry of the article is transcribed and no reference is redistributed. Reference adaptation: none was needed and none was made; every reference inside the delivered unit resolves inside it, so no unresolved reference or citation is produced. Printed slips kept literal: no printed word, symbol, hypothesis or number of the counted statement or of its proof was altered. Layout and class: the article's own class is \texttt{amsart} with packages including amsmath, amssymb, cite, graphics, epsfig, color, euscript, xcolor, enumitem, a4wide, tikz; the wrapper is the lane's pinned \texttt{amsart} editorial wrapper at 10pt on A4, which restates the theorem-like environments the retained text needs and adds the article's own macro definitions that the retained text needs (\texttt{\textbackslash bb}, \texttt{\textbackslash mc}, \texttt{\textbackslash mf}, \texttt{\textbackslash ms}, \texttt{\textbackslash mtt}), with extra wrapper packages (bm, bbm, dsfont, xspace, cancel, mathtools, xcolor). The delivered numbering is the unit's own. Assets: no retained span contains a figure environment, a graphics-inclusion command, a PDF-page inclusion, a table or a TikZ picture; no image file is copied into this package, the package declares no asset dependency and no image file is redistributed with it. Licence: version arXiv:2505.20217v1 of this article is distributed under the Creative Commons Attribution 4.0 International licence, as recorded in the arXiv licence metadata for this exact version and shown on the abstract page https://arxiv.org/abs/2505.20217v1. A full rights notice is delivered at licenses/arxiv-2505.20217-CC-BY-4.0.txt. Characters: every retained excerpt is pure ASCII, so the delivered engine, XeLaTeX with its default Latin Modern text font, reports no missing character.
\begin{equation}
\label{71}
{\color{blue}\ms E(m)} \,:=\, N(m ,r_0)\,\cap\, B(m ,r_0)\;.
\end{equation}

\begin{corollary}
\label{l07}
Fix $1\le p\le \mf q$. Then, for all $k\in\bb Z$,
\begin{equation*}
\lim_{a \to 0} \limsup_{\epsilon \to 0} 
\sup_{x\in \ms E(\ms M_p(k))} 
\bb P^{p,\epsilon}_{x}\big[\, \tau(\ms E_p \setminus \ms E(\ms M_p(k)))
\le  a  \,\big] \,=\, 0\;.
\end{equation*}
\end{corollary}

\begin{proposition}
\label{l11}
Fix $1\le p\le \mf q$, $r_0>0$ as in \eqref{71}, and $k\in \bb Z$. Then,
\begin{equation*}
\lim_{\epsilon\rightarrow0}
\mathbb{P}_{x_\epsilon}^{p,\epsilon}
\Big[\, \bigcap_{j=1}^{{\mf n}}
\big \{\, \mtt x (t_{j,\epsilon} )
\in \mc E( \ms M_p(k_j))\, \big\} \, \Big]
\,=\, \mtt Q^{(p)}_{\ms M_p(k)}
\big[\, \bigcap_{j=1}^{{\mf n}}
\big\{ \, \mtt x (t_{j}) = \ms M_p(k_j) \, \} \, \, \big]
\end{equation*}
for all ${\mf n}\ge 1$, $0<t_{1}<\cdots<t_{{\mf n}}$,
$k_1, \dots, k_{\mf n} \in \bb Z$, and sequences
$x_\epsilon \in \mc E(\ms M_p(k))$, $t_{j,\epsilon} \to t_j$.
\end{proposition}

\begin{equation}
\label{33}
\limsup_{\epsilon\to0}\,
\sup_{t\ge c}\, \sup_{x\in \bb R}\,
\bb P_{x}^{p, \epsilon}    [\, \mtt x(t) \notin\ms E_p \, ] \;=\; 0\;.
\end{equation}

\begin{lemma}
\label{l36}
Fix $1\le p\le \mf q$, $j$, $t>0$, $k\in\bb Z$ and $m \in \ms M_p(k)$.
Then, for any sequence $(\vartheta^{(p)}_{\epsilon} :\epsilon>0)$ such
that $\vartheta^{(p)}_{\epsilon}\prec\theta_{\epsilon}^{(p)}$,
\begin{equation*}
\lim_{\epsilon\to 0}\, \sup_{x \in\mathcal{E}(\ms M_p(j))}\,
\sup_{|s|\le\vartheta^{(p)}_\epsilon}\,\Big|\,
\mathbb{P}_{x}^{\epsilon} \big[\,\mtt x(\theta_{\epsilon}^{(p)}t+s)
\in\mathcal{E}(m)\,\big] \,-\,
\frac{\pi (m)}{\pi(\ms M_p(k))}
\,\mtt {Q}_{\ms M_p(j)}^{(p)}\big[\, \mtt
x(t)= \ms M_p(k)\,\big]\,\Big|\ =\ 0\,. 
\end{equation*}
\end{lemma}

\begin{proof}
The proof is by induction. For $p=1$, as the sets 
$\ms M_1(\ell)$ are singletons, the assertion of the lemma
corresponds to the one of Proposition \ref{l11}.

Fix $p\ge2$ and assume that the lemma holds for $1\le p'<p$.  Fix 
$k\in \bb Z$, $m \in \ms M_p(k)$ and a sequence
$(\vartheta^{(p)}_\epsilon)_{\epsilon>0}$ such that
$\vartheta^{(p)}_\epsilon\prec\theta_{\epsilon}^{(p)}$.  By
construction,
$\ms M_p(k) = \cup_{\ell_1 \le i\le \ell_2} \ms M_{p-1}(i)$ for some
$\ell_1\le\ell_2$. Denote by $i_0 \in \{\ell_1, \dots, \ell_2\}$ the
index such that $m \in \ms M_{p-1}(i_0)$.

Fix $\eta>0$. Since
$\{ \ms M_{p-1}(i) : \ell_1 \le i\le \ell_2\}$ forms a
$\bb X_{p-1}$-recurrent class, by the paragraph preceding the
statement of the lemma, there exists $s_{0}=s_{0}(\eta)>0$ such that
\begin{equation}
\label{93}
\max_{\ell_1 \le i \le \ell_2}\,\Big |\,
\mtt Q_{\ms M_{p-1}(i)}^{(p-1)} \big[\, \mtt x (s_{0}) =
\ms M_{p-1}(i_0)\,\big]
\,-\, \frac{\pi(\ms M_{p-1}(i_0))}{\pi(\ms M_p(k))}\,\Big | \,\le\,
\eta\,.
\end{equation}

Fix $x\in \ms E(\ms M_p(k))$, $|s|\le\vartheta^{(p)}_\epsilon$,
and write
\begin{equation*}
\bb P_x^{\epsilon }\big[\, \mtt x(\theta_{\epsilon}^{(p)}t+s) \in
\ms E(m) \,\big]
\,=\,  \sum_{j=1}^{3} \bb P_x^{\epsilon } \big[\, \mtt x
(\theta_{\epsilon}^{(p)}t+s) \in \ms E(m)
\,\big|\,\mathscr{A}_{j}\,\big]\,
\bb P_x^{\epsilon } [\,\mathscr{A}_{j}\,]\;,
\end{equation*}
where
\begin{gather*}
\mathscr{A}_{1}\ =\ \{\, \mtt x (t_{\epsilon}) \in \ms E(\ms M_{p}(k))
\,\} \;,\quad \mathscr{A}_{2}\ =\ \{\, \mtt x (t_{\epsilon}) \in \ms
E(\ms M_{p}) \setminus \ms E(\ms M_{p}(k)) \,\}\;, \quad
\mathscr{A}_{3}\ =\ \{\, \mtt x (t_{\epsilon})\not\in \ms E(\ms M_{p})
\,\}\;,
\end{gather*}
and $t_{\epsilon}=\theta_{\epsilon}^{(p)}t+s-\theta_{\epsilon}^{(p-1)}s_{0}$.

We claim that the term corresponding to $\mathscr{A}_{2}$ is
negligible.  Bound $\bb P_x^{\epsilon } [\,\mathscr{A}_{2}\,]$ by
$1$ and apply the strong Markov property to get that this term is less
than or equal to
\begin{equation*}
\sup_{k'\neq k}\,\sup_{y\in \ms M_p(k')}\,
\mathbb{P}_{y}^{\epsilon}\big[\,\tau_{\ms E(\ms M_p(k))}
<\theta_{\epsilon}^{(p-1)}s_{0}\,\big]\;.
\end{equation*}
By Corollary \ref{l07} this expression vanishes as $\epsilon\to 0$.

By \eqref{33}, the term corresponding to $\mathscr{A}_{3}$ is also
negligible.  We turn to the term corresponding to
$\mathscr{A}_{1}$. We estimate the two probabilities separately. For
the first probability, by the strong Markov property,
\begin{align*}
& \sup_{|s|\le\vartheta^{(p)}_\epsilon}\, \Big|\,
\bb P_x^{\epsilon } \big[\, \mtt x
(\theta_{\epsilon}^{(p)}t+s) \in \ms E(m)
\,\big|\,\mathscr{A}_1 \,\big]
\,-\, \frac{\pi(m)}{\pi(\ms M_{p}(k))}
\, \Big|
\\
&\quad \le
\sup_{y\in \ms E(\ms M_p(k))}\, \Big|\,
\bb P_y^{\epsilon } \big[\, \mtt x
(\theta_{\epsilon}^{(p-1)}s_0) \in \ms E(m) \,\big]
\,-\, \frac{\pi(m)}{\pi(\ms M_{p}(k))}
\, \Big|\;.
\end{align*}
This expression is less than or equal to
\begin{align*}
& \max_{\ell_1\le i\le \ell_2}\, \sup_{y\in \ms E(\ms M_{p-1}(i))}
\,\Big|\, \bb P_y^{\epsilon } \big[\, \mtt x
(\theta_{\epsilon}^{(p-1)}s_0) \in \ms E(m) \,\big]
\,-\, \frac{\pi(m)}{\pi(\ms M_{p-1}(i_0))} 
\, \mtt Q^{(p-1)}_{\ms M_{p-1}(i)} \big[\,
\mtt x (s_{0}) = \ms M_{p-1}(i_0) \,\big] \,\Big|
\\
& \quad +\, \max_{\ell_1\le i\le \ell_2}\, \sup_{y\in \ms E(\ms M_{p-1}(i))}
\,\Big|\, \frac{\pi(m)}{\pi(\ms M_{p-1}(i_0))} 
\, \mtt Q^{(p-1)}_{\ms M_{p-1}(i)} \big[\,
\mtt x (s_{0}) = \ms M_{p-1}(i_0) \,\big]
\,-\,  \frac{\pi(m)}{\pi(\ms M_{p}(k))} \,\Big|\;.
\end{align*}
By the induction hypothesis, the first term converges to $0$ as
$\epsilon\to0$.  Therefore, by \eqref{93},
\begin{equation*}
\limsup_{\epsilon\to0}  \sup_{x \in\mathcal{E}(\ms M_p(j))}
\sup_{|s|\le\vartheta^{(p)}_\epsilon}\, \Big|\,
\bb P_x^{\epsilon } \big[\, \mtt x
(\theta_{\epsilon}^{(p)}t+s) \in \ms E(m)
\,\big|\,\mathscr{A}_1 \,\big]
\,-\, \frac{\pi(m)}{\pi(\ms M_{p}(k))}
\, \Big| \,\le\, \eta\;.
\end{equation*}

For the second probability, by Proposition \ref{l11},
\begin{equation*}
\limsup_{\epsilon\to0}  \sup_{x \in\mathcal{E}(\ms M_p(j))}
\sup_{|s|\le\vartheta^{(p)}_\epsilon}\, \Big|\,
\bb P_x^{\epsilon } \big[\, \mtt x (t_{\epsilon}) \in \ms E(\ms
M_{p}(k))  \,\big]
\,-\, \mtt Q^{(p-1)}_{\ms M_{p}(j)} \big[\,
\mtt x (s_{0}) = \ms M_{p}(k) \,\big]
\, \Big| \,=\, 0\;.
\end{equation*}
To complete the proof of the induction step and the one of the lemma,
it remains to recollect all the previous estimates, letting $\eta\to 0$
at the end.
\end{proof}

\end{document}
